\part{Loop Models / Anyonic Chains}

\section{Temperley--Lieb}

\begin{exercise}[The Temperley--Lieb diagram category and the diagram basis of $TL_n(d)$]
Fix $d\in\R$. A \emph{planar pairing} from $n$ to $m$ is an isotopy class of $(n+m)/2$ pairwise disjoint smooth arcs in the strip $[0,1]\times[0,1]$ whose endpoints are exactly the points $(i/(n+1),0)$, $1\le i\le n$, and $(j/(m+1),1)$, $1\le j\le m$; it exists only when $n+m$ is even. Define the category $\mathcal{TL}(d)$: objects are the integers $n\ge0$; $\Hom(n,m)$ is the complex vector space with basis the planar pairings from $n$ to $m$; the composite $g\circ f$ of pairings $f\colon n\to m$ and $g\colon m\to p$ is obtained by stacking $g$ on top of $f$, rescaling, deleting each closed loop so formed, and multiplying by $d$ for each deleted loop; composition is extended bilinearly. The tensor product $n\otimes m=n+m$ places diagrams side by side. The involution $f\mapsto f^*$ reflects a pairing in the line $\{t=1/2\}$ and is extended antilinearly. Write $\cap\in\Hom(2,0)$ and $\cup\in\Hom(0,2)$ for the single arc, $1_n$ for the identity pairing, and $e_i=1_{i-1}\otimes(\cup\circ\cap)\otimes1_{n-i-1}\in\End(n)$ for $1\le i<n$.

Prove that composition is associative and unital, that $\otimes$ is a strict monoidal structure with $(g\circ f)\otimes(g'\circ f')=(g\otimes g')\circ(f\otimes f')$, and that $*$ is an antilinear contravariant monoidal involution. Prove the snake identities $(\cap\otimes1_1)\circ(1_1\otimes\cup)=1_1=(1_1\otimes\cap)\circ(\cup\otimes1_1)$ and that $e_i^2=de_i$, $e_ie_{i\pm1}e_i=e_i$, $e_ie_j=e_je_i$ for $\abs{i-j}\ge2$. Deduce a homomorphism $\phi_n\colon TL_n(d)\to\End(n)$ from the algebra $TL_n(d)$ of the section Hecke / Temperley--Lieb algebras in Part V. Prove that every planar pairing $n\to n$ other than $1_n$ is a product of the $e_i$, by induction on the number of arcs joining two bottom points: choose an innermost such arc, joining adjacent points $i,i+1$, and factor the diagram through $e_i$. Prove that the number of planar pairings from $n$ to $n$ is the Catalan number $C_n=\binom{2n}{n}/(n+1)$ by exhibiting a bijection with planar pairings of $2n$ points on a line (rotate the bottom row to the right of the top row) and deriving the recursion $C_{n}=\sum_{k}C_{k-1}C_{n-k}$ from the arc through the first point. Using the dimension $C_n$ of $TL_n(d)$ established in Part V, conclude that $\phi_n$ is an isomorphism and that the planar pairings form a basis of $TL_n(d)$, the \emph{diagram basis}.
\end{exercise}

\begin{exercise}[Jones--Wenzl idempotents]
Let $q\in\C^\times$ with $d=q+q^{-1}$ and $[n]_q=(q^n-q^{-n})/(q-q^{-1})$, so $[2]_q=d$ and $[n+1]_q=d[n]_q-[n-1]_q$. Assume $[k]_q\ne0$ for $1\le k\le n$. Define $f_0=1_0$, $f_1=1_1$ and, for $1\le m<n$,
\[
f_{m+1}=f_m\otimes1_1-\frac{[m]_q}{[m+1]_q}\,(f_m\otimes1_1)\,e_m\,(f_m\otimes1_1)\in TL_{m+1}(d).
\]
Prove by simultaneous induction on $m\le n$ the following statements. (i) $f_m$ is an idempotent with $e_if_m=f_me_i=0$ for $1\le i<m$, and $f_m-1_m$ lies in the two-sided ideal of $TL_m(d)$ generated by $e_1,\dots,e_{m-1}$. (ii) $f_m$ is the unique element with these two properties. (iii) The partial closure identity
\[
(1_{m-1}\otimes\cap)\,(f_m\otimes1_1)\,(1_{m-1}\otimes\cup)=\frac{[m+1]_q}{[m]_q}\,f_{m-1}
\]
holds in $TL_{m-1}(d)$. (iv) $f_m(f_{m-1}\otimes1_1)=f_m=(f_{m-1}\otimes1_1)f_m$ and $f_mTL_m(d)f_m=\C f_m$; in particular $\End(f_m)=\C$ in the additive idempotent completion of $\mathcal{TL}(d)$, whose objects are finite direct sums of pairs $(n,p)$ with $p\in\End(n)$ idempotent and whose morphisms $(n,p)\to(n',p')$ are the elements $p'\Hom(n,n')p$.

For the induction step of (iii) expand $f_{m+1}$ by its definition, use $e_m(f_m\otimes 1_1)e_m=(1_{m-1}\otimes\cup)\,[(1_{m-1}\otimes\cap)(f_m\otimes1_1)(1_{m-1}\otimes\cup)]\,(1_{m-1}\otimes\cap)$ together with the induction hypothesis, and simplify with $[m+2]_q=d[m+1]_q-[m]_q$.

Put $p=\frac{[m]_q}{[m+1]_q}(f_m\otimes1_1)e_m(f_m\otimes1_1)$ for $m<n$, so that $f_m\otimes1_1=f_{m+1}+p$. Prove that $p$ is an idempotent orthogonal to $f_{m+1}$ and that the morphisms $\iota=\frac{[m]_q}{[m+1]_q}(f_m\otimes1_1)(1_{m-1}\otimes\cup)\colon(m-1,f_{m-1})\to(m+1,p)$ and $\pi=(1_{m-1}\otimes\cap)(f_m\otimes1_1)\colon(m+1,p)\to(m-1,f_{m-1})$ satisfy $\pi\iota=f_{m-1}$, $\iota\pi=p$. Conclude the isomorphism $(m,f_m)\otimes(1,f_1)\simeq(m-1,f_{m-1})\oplus(m+1,f_{m+1})$ in the additive idempotent completion, for $1\le m<n$.
\end{exercise}

\begin{exercise}[The Markov trace, negligible morphisms, and the semisimplification at roots of unity]
Define the \emph{Markov trace} $\operatorname{tr}_n\colon TL_n(d)\to\C$ on a planar pairing $x\colon n\to n$ by joining the $i$-th top point to the $i$-th bottom point by $n$ nested arcs to the right of the diagram and setting $\operatorname{tr}_n(x)=d^{\ell(x)}$, $\ell(x)$ the number of closed loops so obtained. Prove that $\operatorname{tr}_n(xy)=\operatorname{tr}_n(yx)$, $\operatorname{tr}_{n+1}(x\otimes1_1)=d\operatorname{tr}_n(x)$, $\operatorname{tr}_{n+1}(x e_n)=\operatorname{tr}_n(x)$ for $x\in TL_n(d)$, $\operatorname{tr}_{n+m}(x\otimes y)=\operatorname{tr}_n(x)\operatorname{tr}_m(y)$, and $\operatorname{tr}_n(x^*)=\overline{\operatorname{tr}_n(x)}$ for real $d$. Prove from the partial closure identity of the preceding exercise that $\operatorname{tr}_m(f_m)=[m+1]_q$ whenever $f_m$ is defined.

Call $x\in\Hom(n,m)$ \emph{negligible} if $\operatorname{tr}_n(yx)=0$ for all $y\in\Hom(m,n)$. Prove that negligible morphisms form a two-sided ideal $\mathcal N$ closed under $\otimes$ with arbitrary morphisms and under $*$ (for the tensor property, write $\operatorname{tr}_{n+p}(z(x\otimes1_p))$ as $\operatorname{tr}_n(z'x)$ with $z'$ a partial closure of $z$), so that the quotient $\mathcal{TL}(d)/\mathcal N$ is again a $\C$-linear monoidal category with involution.

Now let $k\ge1$ and $q=e^{\ii\pi/(k+2)}$, $d=2\cos(\pi/(k+2))$, so that $[m]_q=\sin(m\pi/(k+2))/\sin(\pi/(k+2))>0$ for $1\le m\le k+1$ and $[k+2]_q=0$. Prove that $f_{k+1}$ is negligible: for $y\in TL_{k+1}(d)$ write $f_{k+1}yf_{k+1}=\lambda(y)f_{k+1}$ and compute $\operatorname{tr}(f_{k+1}y)=\lambda(y)[k+2]_q$. Prove that in the additive idempotent completion of $\mathcal{TL}(d)/\mathcal N$: (a) $\Hom((a,f_a),(b,f_b))=0$ for $a\ne b$ in $\{0,\dots,k\}$ (every pairing $a\to b$ with $a\neq b$ has an arc joining two bottom or two top points); (b) $\End((a,f_a))=\C$; (c) every object $(n,1_n)$ is isomorphic to a direct sum of objects $(a,f_a)$, $0\le a\le k$, by induction on $n$ using $(m,f_m)\otimes(1,f_1)\simeq(m-1,f_{m-1})\oplus(m+1,f_{m+1})$ and $(k+1,f_{k+1})\simeq0$. Write $\mathcal C_k$ for this category and $a$ for the object $(a,f_a)$; call the objects $0,\dots,k$ simple and $\mathcal C_k$ semisimple. Prove the fusion rule
\[
a\otimes b\simeq\bigoplus_{\substack{c\equiv a+b\ (\mathrm{mod}\ 2)\\ \abs{a-b}\le c\le\min(a+b,\,2k-a-b)}}c
\]
by induction on $b$, using $a\otimes(b\otimes1)\simeq(a\otimes b)\otimes 1$ and the rule for $\otimes1$. Compare with the fusion rules of the $A_{k+1}$ RSOS heights of the section Quantum dimensions in Part VI, and verify $\operatorname{tr}(f_a)=[a+1]_q$ against the quantum dimensions $d_a=[a+1]_q$ computed there.
\end{exercise}

\begin{exercise}[The dense loop model and its six-vertex representation]
Let $\Lambda$ be the square lattice rotated by $\pi/4$, drawn on the strip $\{1,\dots,L\}\times\{1,\dots,N\}$ so that each vertex has two edges entering from below (south-west, south-east) and two leaving above (north-west, north-east). A \emph{dense loop configuration} chooses at every vertex one of the two non-crossing pairings of its four half-edges: the pairing $\{\mathrm{SW},\mathrm{NW}\},\{\mathrm{SE},\mathrm{NE}\}$ with weight $\rho_1>0$ or the pairing $\{\mathrm{SW},\mathrm{SE}\},\{\mathrm{NW},\mathrm{NE}\}$ with weight $\rho_2>0$. The union of the chosen arcs is a disjoint union of loops and of open arcs ending on the boundary. With free boundary in the horizontal direction and the strands joined across the top and bottom as in the Markov closure, define
\[
Z_{L,N}(n;\rho_1,\rho_2)=\sum_{\text{configurations}}\rho_1^{\#\{\rho_1\text{-vertices}\}}\rho_2^{\#\{\rho_2\text{-vertices}\}}\,n^{\#\{\text{loops}\}} .
\]
Prove that $Z_{L,N}=\operatorname{tr}_L(T^N)$ with $d=n$ and
\[
T=\prod_{i\ \mathrm{even}}(\rho_1+\rho_2e_i)\prod_{i\ \mathrm{odd}}(\rho_1+\rho_2e_i)\in TL_L(n)
\]
(the two products over the two sublattices of vertices of a double row), identifying the vertex pairings with $1$ and $e_i$.

Now let $n=2\cos\gamma$ with $0\le\gamma<\pi$ and orient every loop in one of its two directions. Give each oriented loop the weight $e^{\ii\gamma\,\omega/(2\pi)}$ where $\omega\in\{\pm2\pi\}$ is its total turning angle, and prove that summing over the two orientations of a contractible loop yields $n$. Draw the local pairings with quarter-circle arcs. Prove that each passage through such an arc turns by $\pm\pi/2$, so that the loop weight factorizes as a product over vertices of $e^{\pm\ii\gamma/4}$ per turn. Orienting each edge by its loop, prove that every vertex has two entering and two leaving arrows (the ice rule of Part IV), and that the sum over pairings compatible with a given arrow configuration, weighted by $\rho_i$ and the turn factors, produces the six-vertex weights
\[
a=\rho_1,\qquad b=\rho_2,\qquad c_\pm=\rho_1e^{\pm\ii\gamma/2}+\rho_2e^{\mp\ii\gamma/2},
\]
where $a$ ($b$) is the weight of the two vertices at which both lines pass straight with parallel (antiparallel) orientation and $c_\pm$ are the weights of the two vertices with sources or sinks along a line. Prove that in every configuration on a domain with prescribed boundary arrows the difference between the numbers of $c_+$ and $c_-$ vertices is determined by the boundary, so that $Z$, after division by the boundary factor $(c_+/c_-)^{(N_+-N_-)/2}$, depends on $c_\pm$ only through $c^2:=c_+c_-$, and compute
\[
\Delta=\frac{a^2+b^2-c^2}{2ab}=-\cos\gamma=-\frac n2 .
\]
State precisely which loops receive the weight $n$ under this correspondence (contractible loops) and which receive the weight $2$ (loops whose turning angle vanishes, that is, loops winding around the periodic direction), and conclude the equality of the dense loop partition function on the plane with the six-vertex partition function of Part IV at $\Delta=-n/2$, for $0<n\le2$.
\end{exercise}

\section{Fusion graphs}

\begin{exercise}[Fusion rings, fusion matrices, and fusion graphs]
A \emph{fusion ring} is a ring $K$ that is free as a $\Z$-module on a finite set $B\ni1$ containing the unit, such that the structure constants $ab=\sum_{c\in B}N_{ab}^c\,c$ are nonnegative integers and such that there is an involution $a\mapsto\bar a$ of $B$, extending to an anti-automorphism of $K$, with $N_{ab}^1=\delta_{b,\bar a}$. Elements of $B$ are the simple elements. Prove the reciprocity identities $N_{ab}^c=N_{\bar a c}^b=N_{c\bar b}^a=N_{\bar b\bar a}^{\bar c}$ (write $N_{ab}^c$ as the coefficient of $1$ in $(ab)\bar c$ and reassociate). Define the \emph{fusion matrices} $N_a\in\Mat_B(\Z_{\ge0})$ by $(N_a)_{cb}=N_{ab}^c$, so that $N_a$ is the matrix of left multiplication by $a$ in the basis $B$; prove $N_aN_b=\sum_cN_{ab}^cN_c$ and $N_{\bar a}=N_a^{\mathsf T}$. Define the \emph{fusion graph} $\Gamma_x$ of $x\in B$: vertices $B$ and $N_{xa}^b$ edges from $a$ to $b$; it is undirected when $\bar x=x$. Call $x$ a \emph{generator} if every $b\in B$ occurs in some power $x^m$, and prove that $x$ is a generator if and only if $\Gamma_x$ is strongly connected.

Prove that $K$ is commutative if and only if all $N_a$ commute. For the fusion rule of the section Temperley--Lieb at level $k$ (simples $0,\dots,k$, $\bar a=a$), prove that it defines a commutative fusion ring, that $1$ is a generator, and that $\Gamma_1$ is the path graph with vertices $0,\dots,k$ and edges $\{a,a+1\}$, that is, the Dynkin diagram $A_{k+1}$ of the RSOS heights $a+1\in\{1,\dots,k+1\}$ of Part VI. Prove that the ring $\Z[x]/(U_{k+1}(x/2))$, with $U_m$ the Chebyshev polynomial of the second kind $U_m(\cos\theta)=\sin((m+1)\theta)/\sin\theta$, is isomorphic to this fusion ring via $x\mapsto1$, with $a\mapsto U_a(x/2)$.
\end{exercise}

\begin{exercise}[Anyonic-chain Hilbert spaces and the Temperley--Lieb path representation]
Let $K$ be a commutative fusion ring, $x=\bar x\in B$ a generator with $N_{xa}^b\in\{0,1\}$ for all $a,b$, and $\Gamma=\Gamma_x$. For $a,b\in B$ and $N\ge0$ define the \emph{anyonic-chain Hilbert space}
\[
\mathcal H_N(a,b)=\mathrm{span}_\C\{(a_0,\dots,a_N):a_0=a,\ a_N=b,\ N_{xa_{i-1}}^{a_i}=1\text{ for }1\le i\le N\},
\]
the space of functions on admissible paths of length $N$ in $\Gamma$ from $a$ to $b$, with the paths orthonormal. Prove $\dim\mathcal H_N(a,b)=(N_x^N)_{ba}$ and $\sum_b\dim\mathcal H_N(a,b)=\norm{N_x^N\delta_a}_1$.

Let $\psi\colon B\to(0,\infty)$ satisfy $N_x\psi=\lambda\psi$, that is, $\sum_bN_{xa}^b\psi_b=\lambda\psi_a$ for all $a$. For $1\le i\le N-1$ define $e_i$ on $\mathcal H_N(a,b)$ by
\[
(e_i\Psi)(a_0,\dots,a_N)=\delta_{a_{i-1},a_{i+1}}\sum_{c:\,N_{xa_{i-1}}^c=1}\frac{\sqrt{\psi_{a_i}\psi_c}}{\psi_{a_{i-1}}}\,\Psi(a_0,\dots,a_{i-1},c,a_{i+1},\dots,a_N).
\]
Prove that each $e_i$ is self-adjoint and that $e_i^2=\lambda e_i$, $e_ie_{i\pm1}e_i=e_i$, $e_ie_j=e_je_i$ for $\abs{i-j}\ge2$, so that $\mathcal H_N(a,b)$ is a $*$-representation of $TL_N(\lambda)$, the \emph{path representation}. For the level-$k$ fusion ring with $x=1$ and $\psi_a=\sin(\pi(a+1)/(k+2))$, the sine eigenvector of Part VI with $\lambda=2\cos(\pi/(k+2))$, verify that $e_i/\lambda$ is the orthogonal projector onto the subspace of paths in which the two steps $a_{i-1}\to a_i\to a_{i+1}$ return to $a_{i-1}$ with the intermediate label weighted by $\sqrt{\psi_{a_i}}$; that is, $e_i/\lambda$ is a rank-one projector on each fibre $\{(a_{i-1},\cdot,a_{i+1}):a_{i-1}=a_{i+1}\}$ and zero on the other fibres. Call $P^{(1)}_i=e_i/\lambda$ the projector onto the trivial channel at bond $i$: it is the orthogonal projector onto the states in which the two $x$-labels at positions $i,i+1$ fuse to $1$. Prove $\Tr P^{(1)}_i=\sum_c(N_x^{i-1})_{ca}(N_x^{N-i-1})_{bc}=(N_x^{N-2})_{ba}$ by counting the fibres on which $P^{(1)}_i\ne0$.
\end{exercise}

\begin{exercise}[The golden chain as a Temperley--Lieb Hamiltonian]
The \emph{Fibonacci fusion ring} has $B=\{1,\tau\}$, $\bar\tau=\tau$, and $\tau^2=1+\tau$. Prove it is a fusion ring, that $\tau$ is a generator, that $\Gamma_\tau$ has the single edge $\{1,\tau\}$ and a loop at $\tau$, that $\psi=(1,\varphi)$ with $\varphi=(1+\sqrt5)/2$ is a positive eigenvector of $N_\tau$ with eigenvalue $\varphi$, and that $\dim\mathcal H_N(1,\tau)=F_N$, $\dim\mathcal H_N(1,1)=F_{N-1}$ for $N\ge1$, with $F_0=0$, $F_1=F_2=1$, $F_{N+1}=F_N+F_{N-1}$. Prove that the map $1\mapsto0$, $\tau\mapsto2$ identifies the Fibonacci ring with the subring of the level-$3$ ring of the section Temperley--Lieb spanned by the even simples $0,2$, and that $\varphi=[3]_q$ at $q=e^{\ii\pi/5}$.

The \emph{golden chain} is $H=-\sum_{i=1}^{N-1}P^{(1)}_i$ on $\mathcal H_N(a,b)$ (open chain) or on the space of closed admissible paths $(a_0,\dots,a_N=a_0)$ with $i$ taken modulo $N$ (periodic chain). Prove that $H=-\varphi^{-1}\sum_ie_i$ with $e_i$ the path representation of $TL_N(\varphi)$, and write out the $2\times2$ block of $P^{(1)}_i$ on the fibre $a_{i-1}=a_{i+1}=\tau$, $a_i\in\{1,\tau\}$, and its value on the fibre $a_{i-1}=a_{i+1}=1$:
\[
P^{(1)}_i\big|_{(\tau,\cdot,\tau)}=\begin{pmatrix}\varphi^{-2}&\varphi^{-3/2}\\ \varphi^{-3/2}&\varphi^{-1}\end{pmatrix},\qquad P^{(1)}_i\big|_{(1,\tau,1)}=1,
\]
and $P^{(1)}_i=0$ on paths with $a_{i-1}\ne a_{i+1}$. Verify $\varphi^{-1}+\varphi^{-2}=1$ and the projector property directly. Prove that path reversal intertwines the Hamiltonians on $\mathcal H_N(a,b)$ and $\mathcal H_N(b,a)$, and that on the periodic chain $H$ commutes with the translation $(a_i)\mapsto(a_{i+1})$. Compute the spectrum of $H$ for $N=3$ on $\mathcal H_3(1,\tau)$ and on $\mathcal H_3(\tau,\tau)$.
\end{exercise}

\section{Perron--Frobenius}

\begin{exercise}[The Perron--Frobenius theorem for irreducible nonnegative matrices]
Let $B$ be a finite vertex set of size $n$ and $A:\R^B\to\R^B$ a kernel with nonnegative values. Call $A$ \emph{irreducible} if for all $i,j$ there is $m\ge1$ with $(A^m)_{ij}>0$, and \emph{primitive} if there is $m\ge1$ with all entries of $A^m$ positive. Write $x\ge0$ and $x>0$ for entrywise inequalities on vectors, $\abs y$ for the vector of absolute values, and $r(A)=\max\{\abs\lambda:\lambda\in\spec A\}$. Assume $A$ irreducible.
\begin{enumerate}[label=(\alph*)]
\item Prove that $P=(I+A)^{n-1}$ has all entries positive, and that $A$ has no zero row and no zero column.
\item For $x\ge0$, $x\ne0$, put $f(x)=\min\{(Ax)_i/x_i:x_i>0\}$. Prove $f(tx)=f(x)$ for $t>0$, $f(Px)\ge f(x)$, and that $f$ is continuous on $\{x>0\}$. Prove that $r:=\sup\{f(x):x\ge0,x\ne0\}$ is attained: the supremum over the simplex $\{x\ge0,\sum_ix_i=1\}$ is bounded by the supremum over its image under $P$, a compact subset of $\{x>0\}$.
\item Let $x$ attain $r$. Prove $Ax\ge rx$; if $Ax\ne rx$ then $P(Ax-rx)>0$, hence $A(Px)>r\,Px$ and $f(Px)>r$, a contradiction. Conclude $Ax=rx$, $x=(1+r)^{-(n-1)}Px>0$, and $r>0$.
\item Prove that every eigenvalue $\lambda$ of $A$ satisfies $\abs\lambda\le r$: if $Ay=\lambda y$ then $A\abs y\ge\abs\lambda\abs y$, so $f(\abs y)\ge\abs\lambda$. Prove also that $Ay\ge\rho y$ with $y\ge0$, $y\ne0$ implies $\rho\le r$, and $Ay\le\rho y$ with $y>0$ implies $r\le\rho$ (apply the first part to $A^{\mathsf T}$, which is irreducible, and pair with its positive eigenvector). Hence $r=r(A)$.
\item Prove that the eigenspace of $r$ is one-dimensional: if $Ay=ry$ with $y$ real, choose $t\in\R$ such that $x-ty\ge0$ has a zero entry; then $A(x-ty)=r(x-ty)$, and $x-ty\ne0$ would give $P(x-ty)=(1+r)^{n-1}(x-ty)>0$. Treat complex $y$ through its real and imaginary parts.
\item Prove that $r$ is a simple root of the characteristic polynomial: let $z>0$ satisfy $A^{\mathsf T}z=rz$; if $(A-r)w=x$ for some $w$, then $z^{\mathsf T}x=z^{\mathsf T}(A-r)w=0$, contradicting $z^{\mathsf T}x>0$.
\item Prove that every eigenvector $y\ge0$ of $A$ is a positive multiple of $x$: pair $Ay=\lambda y$ with $z$.
\item Assume $A$ primitive with $A^m>0$. Prove that $\abs\lambda<r$ for every eigenvalue $\lambda\ne r$: if $Ay=\lambda y$ with $\abs\lambda=r$, then $A\abs y\ge r\abs y$ forces $A\abs y=r\abs y$ by (c) applied to $\abs y$, so $\abs y=cx>0$; then $\abs{A^my}=A^m\abs y$ with $A^m>0$ forces all entries of $y$ to have the same argument, so $y\in\C x$ and $\lambda=r$. Deduce $\lim_{m\to\infty}r^{-m}A^m=xz^{\mathsf T}/(z^{\mathsf T}x)$, with an $O(\rho^m)$ error for every $\rho>\abs{\lambda_1}/r$, $\rho<1$, where $\lambda_1$ is a second-largest eigenvalue in modulus; account for possible Jordan blocks.
\item Prove that an irreducible matrix with a positive diagonal entry is primitive, and that for $n>1$ the cyclic permutation kernel of an $n$-cycle is irreducible, not primitive, and has all $n$-th roots of unity as eigenvalues.
\end{enumerate}
\end{exercise}

\begin{exercise}[Frobenius--Perron dimensions of a fusion ring and the growth of anyonic chains]
Let $K$ be a fusion ring with simples $B$, $N_a$ the matrix of left multiplication by $a$, and $R_b$ the matrix of right multiplication by $b$, $(R_b)_{ca}=N_{ab}^c$.
\begin{enumerate}[label=(\alph*)]
\item Prove $N_aR_b=R_bN_a$ for all $a,b$ (associativity), and that $M=\sum_{b\in B}R_b$ has all entries positive: for $c,e\in B$ show $\bar ce\ne0$ by multiplying on the right by $\bar e$, and use $N_{cb}^e=N_{\bar ce}^b$.
\item Let $v>0$ be the Perron--Frobenius eigenvector of $M$. Prove that $N_av$ is a nonnegative nonzero eigenvector of $M$ with eigenvalue $r(M)$, hence $N_av=\lambda_av$ with $\lambda_a>0$, and that $a\mapsto\lambda_a$ extends to a ring homomorphism $\operatorname{FPdim}\colon K\to\R$, the \emph{Frobenius--Perron dimension}, with $\operatorname{FPdim}(1)=1$ and $\operatorname{FPdim}(a)\operatorname{FPdim}(b)=\sum_cN_{ab}^c\operatorname{FPdim}(c)$.
\item Prove that a nonnegative matrix $N$ with a positive eigenvector of eigenvalue $\lambda$ has $r(N)=\lambda$ (bound the entries of $N^m$). Conclude $\operatorname{FPdim}(a)=r(N_a)=r(N_{\bar a})=\operatorname{FPdim}(\bar a)$ and $\operatorname{FPdim}(a)\ge1$ (use $N_{a\bar a}^1=1$).
\item Let $x\in B$ be a generator. Prove that $w=(\operatorname{FPdim}(a))_{a\in B}$ satisfies $N_xw=\operatorname{FPdim}(x)w$ (use $N_{xb}^c=N_{\bar xc}^b$ and multiplicativity), hence is the Perron--Frobenius eigenvector of the irreducible matrix $N_x$ normalized by $w_1=1$. Write $d_a=\operatorname{FPdim}(a)$ and compare it with the quantum dimension of the section Quantum dimensions; for the level-$k$ fusion ring identify $d_a=[a+1]_q$, $q=e^{\ii\pi/(k+2)}$, with the quantum dimensions of Part VI, using the sine eigenvector computed there and (d).
\item Prove that the path representation of the section Fusion graphs exists for every commutative fusion ring and every self-dual generator $x$ with $N_{xa}^b\in\{0,1\}$, with $\psi=w$ and $\lambda=d_x$. Prove $\sum_b\dim\mathcal H_N(a,b)\,d_b=d_x^Nd_a$ exactly, and, when $N_x$ is primitive, $\dim\mathcal H_N(a,b)\sim\dfrac{d_ad_b}{\sum_cd_c^2}\,d_x^{N}$ as $N\to\infty$. Verify for the Fibonacci ring that this reproduces $F_N\sim\varphi^N/\sqrt5$.
\end{enumerate}
\end{exercise}

\begin{exercise}[Nondegenerate leading eigenvalues of transfer matrices and the correlation length]
\leavevmode\par
\begin{enumerate}[label=(\alph*)]
\item Let $T$ be the row-to-row transfer matrix of the six-vertex model of Part IV on $(\C^2)^{\otimes L}$, periodic, with weights $a,b,c>0$, and let $T_M$ be its restriction to the sector with $M$ down arrows, $0\le M\le L$. In the arrow basis, prove that $T_M$ has nonnegative entries and positive diagonal entries. Prove that $T(\sigma',\sigma)>0$ when $\sigma'$ is obtained from $\sigma$ by moving a single down arrow to an adjacent site: the horizontal arrows are determined by the ice rule and alternate at the two flipped sites. Conclude that $T_M$ is irreducible, hence primitive, so that its largest eigenvalue $\Lambda_0(M)$ is positive, simple, and strictly dominant, with positive eigenvector.
\item With $Z=\sum_M\Tr T_M^N$ prove
$\lim_{N\to\infty}N^{-1}\log Z=\log\max_M\Lambda_0(M)$.
For a sector-preserving observable $\mathcal O$, subtract the
leading spectral projection in its connected correlation.
Prove an $O(\rho^r)$ bound for every
$\rho>|\Lambda_1(M)|/\Lambda_0(M)$, $\rho<1$.
When the observable has nonzero overlap with the subleading
spectral subspace, identify the exponential decay rate as
$\log(\Lambda_0(M)/|\Lambda_1(M)|)$; allow polynomial factors
from Jordan blocks and faster decay when that overlap vanishes.
\item Prove that the transfer matrix $T=\prod_{i\ \mathrm{even}}(\rho_1+\rho_2e_i)\prod_{i\ \mathrm{odd}}(\rho_1+\rho_2e_i)$ of the dense loop model, acting in the path representation $\mathcal H_L(a,b)$ of the level-$k$ fusion ring with $\rho_1,\rho_2>0$, has nonnegative entries and positive diagonal in the path basis, and is irreducible on each nonzero fixed-endpoint path space for this path graph and $L\ge2$. Conclude the same spectral statements for the loop model at $n=2\cos(\pi/(k+2))$.
\end{enumerate}
\end{exercise}

\section{Modular tensor categories}

\begin{exercise}[Rigid semisimple categories, ribbon traces, and modularity]
A $\C$-linear category is semisimple here if its objects are finite
direct sums of simple objects, with endomorphism algebra $\C$
for a simple object and no morphisms between nonisomorphic
simples. A monoidal category is rigid if every object has left
and right duals with evaluation and coevaluation satisfying
the snake identities. A pivotal structure is a monoidal
natural isomorphism $j_X:X\to X^{**}$; it is spherical if the
left and right evaluation-coevaluation traces agree.
For example its right trace is
\[
 \operatorname{tr}_X(f)=
 \mathrm{ev}_{X^*}(j_Xf\otimes\id_{X^*})\mathrm{coev}_X.
\]
A ribbon structure consists of a braiding and a compatible
twist $\theta$, with
$\theta_{X\otimes Y}=c_{Y,X}c_{X,Y}(\theta_X\otimes\theta_Y)$
and $\theta_{X^*}=(\theta_X)^*$, where the last star means
the categorical dual morphism.
For simple representatives $a$ let $d_a=\operatorname{tr}_a(\id)$
and choose $\mathcal D$ with $\mathcal D^2=\sum_a d_a^2\ne0$.
Such a category, with finitely many simple classes, simple unit,
finite-dimensional morphism spaces, and invertible kernel
\[
 S_{ab}=\mathcal D^{-1}
          \operatorname{tr}_{a\otimes b}(c_{b,a}c_{a,b}),
\]
is a modular tensor category.
\begin{enumerate}[label=(\alph*)]
\item Prove additivity and multiplicativity of the spherical
dimension by composing evaluation and coevaluation maps.
Recover the loop trace and dimensions of the section
Temperley--Lieb for its rigid diagram category.
\item Let $A=(\Z/2)^2$. On finite-dimensional $A$-graded
spaces use the usual tensor product and duals, and on
homogeneous factors of degrees $(a,b)$ and $(c,d)$ set
$c(v\otimes w)=(-1)^{ad}w\otimes v$ and
$\theta(v)=(-1)^{ab}v$.
Verify the hexagons, balancing, and dual compatibility.
Compute all simple dimensions as one.
\item Show that $\mathcal D=2$ and
$S_{(a,b),(c,d)}=\frac12(-1)^{ad+bc}$.
Prove $S^2=\id$ by summing the four group characters.
\end{enumerate}
\end{exercise}

\begin{exercise}[Fusion diagonalization by a modular kernel in the pointed model]
Use the four-charge category just constructed and let
$(N_xf)(a)=f(a-x)$ on the charge-labelled function space.
\begin{enumerate}[label=(\alph*)]
\item Show that its fusion is group addition, all quantum
dimensions are one, and the columns of $S$ are its
character functions. Prove the kernel identity
$N_xS_{\cdot b}=(S_{xb}/S_{0b})S_{\cdot b}$.
\item Deduce the Verlinde identity in this model,
\[
 N_{xy}^{z}=\sum_b\frac{S_{xb}S_{yb}\overline{S_{zb}}}{S_{0b}},
\]
and evaluate it as $\mathbf1_{x+y=z}$.
\end{enumerate}
\end{exercise}

\section{Braid group representations}

\begin{exercise}[Braiding on fusion spaces and the loop representation]
Use the braid group of the section Hecke / Temperley--Lieb
algebras. For an object $x$ in a braided rigid category,
its fusion space with total charge $a$ is
$\Hom(a,x^{\otimes N})$.
\begin{enumerate}[label=(\alph*)]
\item Construct the braid action by postcomposition with
the adjacent braidings, inserting associators, and verify
the relations using the hexagons.
\item In the Temperley--Lieb diagram category choose
$A\in\C^\times$ with $d=-A^2-A^{-2}$. Prove that
$b_i=A\id+A^{-1}e_i$ satisfies the braid relations and that
$b_i^{-1}=A^{-1}\id+Ae_i$.
If $|A|=1$, $d\in\R$, and $e_i^*=e_i$ in the path
representation, prove that $b_i$ is unitary.
\item In the four-charge category calculate the double
exchange of $e=(1,0)$ and $m=(0,1)$ and obtain $-1$.
Calculate the self-twists of $e$, $m$, and $e+m$.
\end{enumerate}
\end{exercise}

\section{Knot invariants}

\begin{exercise}[The bracket state sum and the Jones normalization]
For an unoriented link diagram $D$ with crossings, a smoothing
chooses one of the two planar pairings at each crossing.
Give the two choices weights $A$ and $A^{-1}$ according to
the crossing convention $b=A\id+A^{-1}e$. If a smoothing
has $\ell$ circles define
\[
 \langle D\rangle=\sum_{\text{smoothings}}
          A^{n_A-n_{A^{-1}}}d^{\ell-1},
 \qquad d=-A^2-A^{-2}.
\]
For an oriented diagram let $w(D)$ be its signed crossing number.
\begin{enumerate}[label=(\alph*)]
\item Prove invariance of the bracket under the second and
third Reidemeister moves using the inverse and braid identities.
Compute the two curl factors as $-A^3$ and $-A^{-3}$.
\item Deduce that
$V_D(A)=(-A^3)^{-w(D)}\langle D\rangle$
is invariant under all three moves, with the positive-crossing
convention chosen to give the curl factor $-A^3$.
Verify that an unknot has value one and a two-component
unlink has value $d$. Define the usual variable by $t=A^{-4}$.
\item Assume $d\ne0$. For the Hopf link obtained by closing two positive
two-strand crossings, calculate
$\operatorname{tr}_2(b_1^2)/d$ and its writhe correction.
Identify the two contributions from the spectral projections
of $e_1$.
\end{enumerate}
\end{exercise}

\begin{exercise}[Prediction: anyonic entropy and a mutual braiding phase]
For the Fibonacci path chain of the section Fusion graphs let
all admissible paths from $1$ to $\tau$ have equal energy zero.
For an independent braiding experiment use the four-charge
category of the section Modular tensor categories, with one
probe of charge $e$ and one enclosed charge $m$.
\begin{enumerate}[label=(\alph*)]
\item Use the fusion recurrence and Perron--Frobenius result
to derive
$\dim\mathcal H_N(1,\tau)=F_N$ and residual entropy per
anyon $k_B\log\varphi$, $\varphi=(1+\sqrt5)/2$.
Determine the leading finite-$N$ correction from
$F_N=(\varphi^N-(-\varphi^{-1})^N)/\sqrt5$.
\item Let a balanced two-path interferometer have reference
relative phase $\phi$ and coherent amplitudes of equal
magnitude. Derive
$P_0(\phi)=(1+\cos\phi)/2$ without an enclosed charge and
$P_m(\phi)=(1-\cos\phi)/2$ when the probe winds once around $m$.
Obtain the phase change from double braiding, equivalently
from $S_{em}S_{00}/(S_{e0}S_{0m})$.
\item Specify independently the assumptions of degenerate
fusion states for the entropy prediction and preserved
charge, adiabatic transport, and equal coherent path amplitudes
for the interference prediction. Calculate the change under
two windings and distinguish it from a change in dynamical phase.
\end{enumerate}
\end{exercise}
